\section{Integration with the BX3 Framework}
\label{sec:bx3_integration}

The BX3 Framework furnishes the structural substrate on which Recursive Spawning operates as a first-class, accountable delegation primitive. Three BX3 layers serve as the backbone of every spawn chain: the Purpose Layer, which declares and anchors intent; the Bounds Engine, which enforces structural and operational constraints at every spawn event; and the Fact Layer, which records every meaningful operation in a tamper-evident forensic ledger. Together they convert a spawn chain from a flexible delegation heuristic into a verifiable accountability infrastructure.

\subsection{Purpose Layer as Accountability Anchor}
\label{sec:purpose_anchor}

The Purpose Layer is the immutable intent anchor of any spawn chain. Every root agent is initialized with a purpose clause—its declared operational scope—that is encoded as the agent's Purpose Layer artifact at provisioning time. This clause is non-negotiable: no child agent may operate outside the semantic scope defined by its parent's purpose, and no agent may revise its own purpose post-spawn.

The Purpose Layer enforces four invariants in Recursive Spawning:

\begin{enumerate}
    \item \textbf{Intent stability.} The purpose clause is frozen at spawn time and remains readable but immutable throughout the agent's lifetime.
    \item \textbf{Downward propagation.} Each child inherits a bounded projection of the parent's purpose clause—computed by the Bounds Engine—not a copy. The projection is deterministic and auditable: any external auditor can verify that the child's scope is a proper subset of the parent's.
    \item \textbf{Upward accountability.} Each agent's Fact Layer record carries a cryptographic reference to its parent's purpose clause, establishing a provable lineage of intent preservation across the full chain.
    \item \textbf{Drift-triggered termination.} An agent whose Fact Layer record exhibits a purpose-layer violation is flagged for deterministic termination. This is not a heuristic—it is a structural property of the chain.
\end{enumerate}

The Purpose Layer thus plays a dual role: it is the source of legitimate operational scope for each agent, and it is the final arbiter of whether any agent in the chain has deviated. No drift can occur without violating the purpose clause, which the Fact Layer makes detectable at the exact point of divergence.

\subsection{Bounds Engine Depth Enforcement}
\label{sec:bounds_depth}

The Bounds Engine enforces the structural and computational constraints that prevent unbounded spawning. In the Recursive Spawning model, the Bounds Engine operates along two axes at every spawn event:

\begin{enumerate}
    \item \textbf{Depth limit.} The maximum chain length $D_{\max}$ is set at root-spawn time and enforced as a precondition for every spawn operation. When an agent attempts to spawn a child at depth $d = D_{\max}$, the Bounds Engine rejects the operation and logs the rejection event in the Fact Layer. This prevents resource exhaustion cascades and ensures that chain traversal always terminates.
    \item \textbf{Scope constraint propagation.} At each spawn event, the Bounds Engine computes the child's operational scope as a constrained projection of the parent's current scope. The projection is deterministic and auditable: given the parent's scope and the spawn operation, any external auditor can verify that $S_{\text{child}} \subseteq S_{\text{parent}}$. The Bounds Engine's scope contract is stored as a Fact Layer artifact at the root, making it itself an auditable artifact.
\end{enumerate}

The Bounds Engine is not a passive filter—it is an active participant in every spawn event. It evaluates the legality of each operation against the current chain state and the root's purpose clause, and it generates a signed constraint artifact that is attached to the child agent's initialization record. This artifact travels with the child throughout its lifetime and is checked at every handoff point in the chain.

A critical property is that the Bounds Engine's scope enforcement is monotonic with respect to chain depth: each generation's scope is a proper subset of its parent's, so the total authorized space monotonically shrinks as the chain grows. This monotonicity is what allows the system to make a definitive statement about any agent's scope relative to the root's original purpose clause.

\subsection{Fact Layer as Forensic Ledger}
\label{sec:fact_forensic}

The Fact Layer is the append-only, tamper-evident ledger that records every meaningful event in a spawn chain. It is the foundation of all correctness guarantees in Recursive Spawning, and its role is structural rather than incidental.

Every Fact Layer entry contains the following fields:

\begin{itemize}
    \item \texttt{event\_type}: The classification of the event (spawn, scope-violation-attempt, termination, purpose-verification, etc.).
    \item \texttt{timestamp}: A globally consistent logical clock value.
    \item \texttt{agent\_id}: The identifier of the agent that generated the event.
    \item \texttt{parent\_pointer}: A cryptographic hash reference to the parent agent's Fact Layer entry at spawn time.
    \item \texttt{purpose\_hash}: The SHA-256 hash of the agent's purpose clause at spawn time.
    \item \texttt{scope\_snapshot}: A serialized snapshot of the agent's operational scope at the time of the event.
    \item \texttt{signature}: The agent's signing key signature over the event payload.
\end{itemize}

The Fact Layer operates as a Merkle-linked log: each entry contains the hash of the previous entry, forming an immutable chain of events. This design ensures:

\begin{enumerate}
    \item \textbf{Non-repudiation.} No agent can deny an action recorded in the Fact Layer, because the entry carries its cryptographic signature.
    \item \textbf{Tamper evidence.} Any modification to a past entry changes its hash, breaking the chain and making the tampering event itself visible.
    \item \textbf{Forensic auditability.} An external auditor can replay the entire chain of events by traversing the Fact Layer from the root to any descendant, verifying each entry's signature and hash-chain integrity.
\end{enumerate}

In the Recursive Spawning context, the Fact Layer records every spawn event with the full provenance of the child agent—including the parent's purpose clause at spawn time, the Bounds Engine's scope constraint, and the child's initialized Purpose Layer artifact. This means any agent that subsequently exhibits drift behavior can be traced to the exact point of divergence by replaying the Fact Layer from the root. There is no gap in the record: the chain of custody is continuous from the root human principal to any active descendant.

The integration of these three BX3 layers creates a structurally coherent system where intent is declared at the top, enforced at every level, and recorded at every step. The Result Layer, when present, consumes Fact Layer records to produce observable outputs—but the Fact Layer itself is the source of truth from which all correctness guarantees are derived.

% ============================================================
\section{Correctness Properties}
\label{sec:correctness}

We establish three correctness properties for any spawn chain constructed under the Recursive Spawning model and integrated with the BX3 Framework as described in Section~\ref{sec:bx3_integration}. We assume \emph{honest initialization}: the root agent's purpose clause is well-formed, the Bounds Engine's depth limit $D_{\max}$ is set to a finite positive integer, and the cryptographic primitives (collision-resistant hash functions, existentially unforgeable signatures) are secure. We do not assume honest behavior of agents post-initialization.

\subsection{Drift Prevention}
\label{sec:drift_prevention}

\begin{prop}[Drift Prevention]
\label{prop:drift_prevention}
Let $C$ be a spawn chain initialized from a root agent $a_0$ with purpose clause $p_0$. For any agent $a_i$ in $C$, if the Fact Layer records of $a_i$ and all of its ancestors are intact—hash-chain verified and signatures valid—then the operational scope of $a_i$ is a subset of $p_0$.
\end{prop}

\begin{Proof}[Proof Sketch]
The proof proceeds by induction on chain depth.

\textbf{Base case ($i = 0$).} $a_0$'s operational scope is defined by $p_0$ by construction, so the property holds trivially.

\textbf{Inductive step.} Assume that for some agent $a_i$ at depth $i$, the operational scope $S_i \subseteq p_0$ holds. When $a_i$ spawns a child $a_{i+1}$, the Bounds Engine computes a constrained projection $S_{i+1} \subseteq S_i$ as the child's scope. By the inductive hypothesis, $S_i \subseteq p_0$, and since $S_{i+1} \subseteq S_i$, it follows that $S_{i+1} \subseteq p_0$. The Fact Layer records the spawn event with the cryptographic parent pointer and the purpose hash of $a_i$, ensuring that any subsequent attempt to operate outside $S_{i+1}$ is logged as a scope-violation-attempt event with a reference to the purpose clause it violates.

Three structural guarantees act in concert to prevent drift:

\begin{enumerate}
    \item \textbf{Immutable parent pointer.} Every spawn event records the parent's Fact Layer entry hash at spawn time. An agent cannot silently re-parent or sever the chain without breaking the Fact Layer hash chain, which is detectable by any external auditor.
    \item \textbf{Forensic ledger completeness.} Every scope-violation attempt is recorded. An agent cannot drift without leaving evidence in the Fact Layer—either a logged violation event or a broken hash chain indicating tampering. Both are detectable.
    \item \textbf{Purpose Layer termination.} The chain terminates at an agent whose Purpose Layer artifact is consistent with $p_0$, because the Bounds Engine rejects any spawn operation that would produce a child with a Purpose Layer inconsistent with its parent's purpose projection. There is no spawn edge that bypasses the Purpose Layer.
\end{enumerate}

Given these three guarantees, for any agent in an intact chain, drift from the root's purpose clause is structurally impossible. $\square$
\end{Proof}

The drift prevention property does not require that agents be trustworthy. It requires only that the cryptographic primitives be secure and that the Fact Layer hash chain be intact. Even a fully compromised agent cannot produce a Fact Layer entry that passes verification without its signing key, which is itself bound to the agent's Purpose Layer artifact at initialization.

\subsection{Chain Integrity}
\label{sec:chain_integrity}

\begin{prop}[Chain Integrity]
\label{prop:chain_integrity}
If a spawn chain $C$ passes the Fact Layer hash-chain verification for all of its entries, then $C$ is bitwise identical to the chain that was constructed during execution. Any modification to any Fact Layer entry, any replay of a past event, or any injection of a false entry will cause hash-chain verification to fail.
\end{prop}

\begin{Proof}[Proof Sketch]
The Fact Layer uses a Merkle-linked log structure where each entry $e_j$ contains $\texttt{prev\_hash} = \text{hash}(e_{j-1})$ alongside its own payload. Verification proceeds by recomputing each hash in sequence and comparing against the stored \texttt{prev\_hash} field. A single mismatched entry invalidates the entire chain from that point onward.

The addition of cryptographic signatures on each entry adds non-repudiation: even if an adversary gains access to the Fact Layer storage, they cannot forge an entry that appears to come from a legitimate agent, because the signature would not verify against the agent's signing key—which is tracked in the agent's Purpose Layer artifact at initialization.

Chain integrity does not depend on the honest behavior of agents. It depends on the integrity of the cryptographic primitives and the hash-chain construction. Even a fully compromised agent cannot produce a Fact Layer entry that passes verification without access to its signing key, which is immutable post-initialization. $\square$
\end{Proof}

This property is critical for external auditors and for downstream consumers of Fact Layer data. A system consuming Fact Layer records can trust them precisely to the extent that the hash-chain verification passes—no further.

\subsection{Termination}
\label{sec:termination}

\begin{prop}[Termination]
\label{prop:termination}
For any spawn chain $C$ with Bounds Engine depth limit $D_{\max}$, the length of $C$ is at most $D_{\max}$. All paths in the chain are finite, and the chain cannot contain cycles.
\end{prop}

\begin{Proof}[Proof Sketch]
The Bounds Engine enforces the depth limit as a strict precondition for every spawn event. A spawn operation at depth $d = D_{\max}$ is atomically rejected: the operation returns an error, and the rejection event is logged in the Fact Layer with the attempting agent's signature. Therefore, no agent at depth $D_{\max}$ can produce a child, establishing the finite upper bound on chain length.

To establish the absence of cycles: each spawn event produces a child with a unique identifier that is bound to the spawner's Fact Layer entry via a cryptographic hash pointer. Since the hash of the parent's Fact Layer entry is stored in the child's initialization record, and since the Fact Layer is append-only and timestamps are strictly increasing, the parent-child relationship is immutable and acyclic by construction. A cycle would require an agent to appear as its own ancestor, which is impossible because each spawn event is logged with a strictly increasing timestamp and an immutable parent pointer—contradicting the well-foundedness of the hash chain. $\square$
\end{Proof}

Termination is a structural guarantee, not a heuristic. It holds regardless of agent behavior, network conditions, or external interference, provided the cryptographic primitives remain secure. The combination of a finite depth limit and an acyclic parent pointer chain ensures that every path in the spawn graph has bounded length.

Together, Propositions~\ref{prop:drift_prevention}, \ref{prop:chain_integrity}, and~\ref{prop:termination} establish that a Recursive Spawning chain governed by the BX3 Framework is a provably accountable, integrity-preserving, and finite lineage of agents—suitable for deployment in high-stakes operational contexts where auditability and correctness are non-negotiable requirements.